\section{Qdrant}

\begin{frame}{Qdrant: Rust-powered Performance Engine}
	\begin{columns}
		\column{0.5\textwidth}
		\begin{block}{Core architecture}
			\begin{itemize}
				\item Single binary với REST API và gRPC API.
				\item Segment là đơn vị lưu trữ tự trị: vector storage, payload storage, HNSW index.
				\item WAL đảm bảo durability trước khi commit.
				\item Rust runtime: memory-safe, không có GC pause.
			\end{itemize}
		\end{block}
		\column{0.5\textwidth}
		\begin{figure}
			\centering
			\fbox{\begin{minipage}{0.9\linewidth}
				\centering
				Qdrant Server\\
				REST / gRPC\\
				Query Planner\\
				Segment 1 \quad Segment 2 \quad Segment N\\
				HNSW + Payload Index\\
				Memmap / On-Disk + WAL
			\end{minipage}}
			\caption{Single-binary architecture của Qdrant.}
		\end{figure}
	\end{columns}
	\note{Qdrant được thiết kế theo hướng gọn, ít dependency và tối ưu latency. Segment giúp chia nhỏ dữ liệu để quản lý index, payload và storage độc lập. WAL giúp dữ liệu được ghi bền vững, còn Rust giảm rủi ro memory bug và tránh pause do garbage collection.}
\end{frame}

\begin{frame}{Filterable HNSW \& ACORN Index}
	\begin{columns}
		\column{0.52\textwidth}
		\begin{block}{Ý tưởng chính}
			\begin{itemize}
				\item HNSW là graph-based ANN index.
				\item Qdrant kiểm tra \textbf{payload filter} trong quá trình graph traversal.
				\item Query Planner chọn strategy theo filter cardinality.
				\item ACORN cải thiện traversal khi filter rất selective.
			\end{itemize}
		\end{block}
		\begin{exampleblock}{Adaptive behavior}
			Filter match nhiều points $\rightarrow$ Filterable HNSW.\\
			Filter match ít points $\rightarrow$ payload index + brute-force nhanh hơn.
		\end{exampleblock}
		\column{0.48\textwidth}
		\begin{figure}
			\centering
			\fbox{\begin{minipage}{0.88\linewidth}
				\centering
				HNSW graph\\
				Matched nodes: traversable\\
				Unmatched nodes: skipped\\
				Connectivity vẫn được tận dụng
			\end{minipage}}
			\caption{Filter được đưa vào quá trình traversal.}
		\end{figure}
	\end{columns}
	\note{Điểm khác biệt của Qdrant không chỉ là có filter, mà là filter được tích hợp vào logic tìm kiếm. Điều này quan trọng với RAG vì metadata như tenant, document type, course id hoặc timestamp thường xuất hiện trong mọi query production.}
\end{frame}

\begin{frame}{Pre-filtering vs Post-filtering vs In-graph Filtering}
	\begin{table}
		\centering
		\caption{Ba chiến lược kết hợp Vector Search với Metadata Filter.}
		\resizebox{\textwidth}{!}{%
		\begin{tabular}{p{0.21\textwidth}p{0.34\textwidth}p{0.33\textwidth}}
			\hline
			\textbf{Strategy} & \textbf{Cách hoạt động} & \textbf{Ảnh hưởng kỹ thuật} \\
			\hline
			Pre-filtering & Lọc metadata trước rồi search trên tập con & Nhanh nếu tập con nhỏ, nhưng HNSW graph có thể bị phân mảnh, làm giảm recall. \\
			Post-filtering & Search vector trước rồi loại kết quả không khớp filter & Dễ triển khai, nhưng có thể thiếu Top-K hợp lệ nếu filter loại nhiều candidates. \\
			In-graph Filtering & Kiểm tra filter trong quá trình duyệt HNSW graph & Giữ graph connectivity tốt hơn, phù hợp filter-heavy RAG workload. \\
			\hline
		\end{tabular}}
	\end{table}
	\begin{alertblock}{Implication for RAG}
		Filtered Search không phải tính năng phụ; nó ảnh hưởng trực tiếp đến Recall@K, latency và tenant isolation.
	\end{alertblock}
	\note{Pre-filtering có thể làm graph bị thưa, post-filtering có thể trả về không đủ candidates hợp lệ. In-graph filtering giữ được lợi thế của graph traversal trong khi vẫn tôn trọng metadata constraints. Đây là lý do Qdrant thường phù hợp workload nhiều filter.}
\end{frame}
